\providecommand{\toprule}{\hline\hline}
\providecommand{\midrule}{\hline}
\providecommand{\bottomrule}{\hline\hline}
\providecommand{\addlinespace}[1][]{}
{\small
\renewcommand{\arraystretch}{1.12}
\begin{tabular}{lrrrr}
\toprule
Ano & G1: Alta (\%) & G2: Média (\%) & G3: Baixa (\%) & Dif. G1 - G3 (p.p.) \\
\midrule
2009 & 33,2 & 32,5 & 34,3 & -1,0 \\
2010 & 33,3 & 32,9 & 33,8 & -0,5 \\
2011 & 33,5 & 33,0 & 33,5 & +0,0 \\
2012 & 33,5 & 33,1 & 33,4 & +0,2 \\
2013 & 33,4 & 32,9 & 33,6 & -0,2 \\
2014 & 33,6 & 32,5 & 33,9 & -0,4 \\
2015 & 33,2 & 31,7 & 35,1 & -1,9 \\
2016 & 33,1 & 31,0 & 35,9 & -2,9 \\
2017 & 33,1 & 30,5 & 36,3 & -3,2 \\
2018 & 33,0 & 30,5 & 36,4 & -3,4 \\
2019 & 33,2 & 30,5 & 36,4 & -3,2 \\
2020 & 33,6 & 30,7 & 35,6 & -2,0 \\
2021 & 33,9 & 30,8 & 35,3 & -1,4 \\
2022 & 34,1 & 31,1 & 34,9 & -0,8 \\
2023 & 34,1 & 31,3 & 34,7 & -0,6 \\
\bottomrule
\end{tabular}
\vspace{2pt}
\par{\footnotesize \textit{Nota:} Participação de cada grupo no emprego formal privado (\%). A classificação de cada ocupação (CBO) é fixada nos tercis de emprego de 2010 (33,3\% em cada grupo em 2010), de modo que a variação no tempo reflete exclusivamente o crescimento relativo das ocupações de cada grupo.}
}
